% Part: normal-modal-logic
% Chapter: filtrations
% Section: introduction

\documentclass[../../../include/open-logic-section]{subfiles}

\begin{document}

\olfileid{nml}{fil}{int}

\olsection{પરિચય}

કોઈ તર્ક વિશેનો એક અગત્યનો પ્રશ્ન હંમેશાં એ છે કે તે નિર્ણેય છે કે
નહિ, એટલે કે ``આ !!{formula} માન્ય છે?'' એ પ્રશ્નનો ઉત્તર આપતી
કોઈ અસરકારક પ્રક્રિયા છે કે નહિ. વિધાનાત્મક તર્ક નિર્ણેય છે:
સત્યાર્થતા સારણી રચીને !!a{formula} પુનરુક્તિ છે કે નહિ તે
અસરકારક રીતે ચકાસી શકાય છે, અને આપેલા !!{formula} માટે સારણી
સાન્ત છે. પરંતુ મોડલ સૂત્ર બધાં નિદર્શનોમાં સત્ય છે કે નહિ તે સીધું
ચકાસી શકાતું નથી, કારણ કે નિદર્શનો અનંત સંખ્યામાં છે. આપેલા સૂત્રને
સંબંધિત બધાં સાન્ત નિદર્શનોની યાદી બનાવી શકાય છે, કારણ કે માત્ર
!!{formula}માં ખરેખર આવતાં !!{propositional variable}sને જગતોના
ઉપગણો નિયુક્ત કરવાથી ફરક પડે છે. પ્રાપ્યતા સંબંધ નિશ્ચિત હોય, તો
$V(p)$ની જુદી જુદી શક્ય નિયુક્તિઓ એ~$W$ના બધા ઉપગણો જ છે;
$\card{W} = n$ હોય તો એવા $2^n$ ઉપગણો છે. આપણા
!!{formula}~$!A$માં $m$ !!{propositional
  variable}s હોય તો~$n$ જગતવાળાં $2^{nm}$ જુદાં નિદર્શનો મળે છે.
દરેક નિદર્શન માટે દરેક જગતમાં~$!A$નું સત્યમૂલ્ય ગણીને સૂત્ર~$!A$
બધાં જગતોમાં સત્ય છે કે નહિ તે ચકાસી શકાય છે. અલબત્ત, બધા શક્ય
પ્રાપ્યતા સંબંધો પણ ચકાસવા પડે; પરંતુ~$n$ જગતો પરના સંબંધો પણ
સાન્ત સંખ્યામાં જ છે (ચોક્કસપણે, $W \times W$ના ઉપગણોની સંખ્યા,
એટલે કે $2^{n^2}$).
\sourcecorrection{OLMD-040}{મૂળના આ ગણતરી-વાક્યમાં ખુલ્લા કૌંસનું
બંધ કૌંસ છૂટી ગયું છે; અહીં તે પુનઃસ્થાપિત કર્યું છે.}

જો આપણને તર્ક~\Log{K} નહિ, પરંતુ નિદર્શનોના કોઈ વર્ગથી
વ્યાખ્યાયિત તર્કમાં રસ હોય (દાખલા તરીકે, સ્વપ્રતિબિંબિત પરંપરિત
નિદર્શનોમાં), તો પ્રાપ્યતા સંબંધ યોગ્ય પ્રકારનો છે કે નહિ તે પણ
ચકાસી શકવું જોઈએ. આપણને રસ ધરાવતી ફ્રેમો મોડલ !!{formula}s વડે
વ્યાખ્યાયિત થઈ શકે ત્યારે એ કરી શકીએ છીએ (દાખલા તરીકે,
\Ax{T} અને \Ax{4} ફ્રેમમાં માન્ય છે કે નહિ તે ચકાસીને). તેથી વિચાર
એવો છે કે બધી સાન્ત ફ્રેમો ક્રમસર લઈએ, દરેક ફ્રેમ આપણને રસ ધરાવતા
વર્ગમાં છે કે નહિ તે ચકાસીએ, પછી તે ફ્રેમ પરનાં બધાં શક્ય
નિદર્શનોની યાદી બનાવી દરેકમાં $!A$ સત્ય છે કે નહિ તે તપાસીએ.
કોઈમાં સત્ય ન હોય તો અટકીએ: $!A$ રસના નિદર્શનોના વર્ગમાં માન્ય
નથી.

આ વિચારમાં સમસ્યા છે: શોધ ક્યારે, કે કદી, બંધ કરવી તે આપણે જાણતા
નથી. જો !!{formula}નું સાન્ત વિરોધી નિદર્શન હોય, તો આપણી પ્રક્રિયા
તે શોધી કાઢશે. પરંતુ સાન્ત વિરોધી નિદર્શન ન હોય તો ઉત્તર નહિ મળે.
!!{formula} માન્ય હોઈ શકે (એકેય વિરોધી નિદર્શન ન હોય), અથવા તેનું
માત્ર અનંત વિરોધી નિદર્શન હોઈ શકે, જે આપણે કદી તપાસવાના નથી.
જો આપણે બતાવી શકીએ કે વિરોધી નિદર્શન ધરાવતું દરેક !!{formula}
સાન્ત વિરોધી નિદર્શન પણ ધરાવે છે, તો આ સમસ્યા દૂર થાય. આવી
પરિસ્થિતિમાં તર્કમાં \emph{સાન્ત નિદર્શન ગુણધર્મ} છે એમ કહીએ છીએ.

પણ તર્કમાં સાન્ત નિદર્શન ગુણધર્મ છે તે કેવી રીતે બતાવવું? એક રસ્તો
એ છે કે~$!A$નું અનંત (વિરોધી) નિદર્શન સાન્ત નિદર્શનમાં ફેરવવાની
રીત શોધવી. એવું કરી શકાય તો જ્યાં $!A$ સત્ય ન હોય એવા કોઈ
નિદર્શનમાંથી મળેલું સાન્ત નિદર્શન પણ $!A$ને અસત્ય જ બનાવશે.
એ સાન્ત નિદર્શન બધાં સાન્ત નિદર્શનોની આપણી યાદીમાં આવશે, અને
આખરે દરેક અમાન્ય સૂત્ર માટે તેની અમાન્યતા નક્કી કરી શકીશું.
જો સૂત્ર માન્ય હોય, તો આ પ્રક્રિયા અટકશે નહિ. જો વધારામાં બતાવી
શકીએ કે પ્રક્રિયાથી મળતા સાન્ત નિદર્શનના કદની કોઈ મહત્તમ મર્યાદા
છે અને તે મર્યાદા માત્ર !!{formula}~$!A$ પર આધારિત છે, તો
વિરોધી નિદર્શન શોધતી વખતે સાન્ત નિદર્શનો કયા કદ સુધી ચકાસવાં
તે જાણીશું. ત્યાં સુધીમાં વિરોધી નિદર્શન ન મળે તો એવું કોઈ
નથી. આમ આપણી પ્રક્રિયા દરેક !!{formula}~$!A$ માટે ``$!A$ માન્ય
છે?'' એ પ્રશ્નનો ખરેખર નિર્ણય કરશે.

અનંત રચનાઓને સાન્ત રચનાઓમાં ફેરવવાની એક ઘણી વાર સફળ થતી
યુક્તિ એ છે કે સંબંધિત બાબતોમાં એકસરખું વર્તન કરતા રચનાના
!!{element}sને એક તરીકે ``ઓળખી'' લઈએ. જો~$\mModel{M}$માં
સંબંધિત બાબતોમાં એકસરખું વર્તન કરતાં અનંત જગતો હોય, તો કદાચ
\emph{બધાં} જગતોને એવા જગતોના સાન્ત સંખ્યામાં (સંભવતઃ અનંત)
``વર્ગો''માં મૂકી શકીએ. બીજા શબ્દોમાં, જગતોના ગણનું યોગ્ય
વિભાજન કરવું જોઈએ: દરેક વર્ગમાં સાન્ત અથવા અનંત જગતો હોઈ શકે,
પણ વર્ગોની સંખ્યા સાન્ત જ હોય.
\sourcecorrection{OLMD-041}{મૂળમાં દરેક વિભાજન-વર્ગમાં અનંત
જગતો હોવા જોઈએ એવું લખાયું છે; સાન્ત વર્ગો પણ માન્ય છે અને અહીં
માત્ર વર્ગોની સંખ્યા સાન્ત હોવી જરૂરી છે.}
પછી નવું નિદર્શન~$\mModel{M^*}$ વ્યાખ્યાયિત કરીએ જેમાં જગતો
આ વર્ગો હોય. જૂના નિદર્શનના સાન્ત સંખ્યામાં વર્ગો નવા નિદર્શનમાં
સાન્ત સંખ્યામાં જગતો આપે છે, એટલે સાન્ત નિદર્શન મળે છે.
જગત~$w$ જે વર્ગમાં હોય તેને $[w]$ કહીએ. આપણે ઇચ્છીએ છીએ કે
$\mSat{M}{!A}[w]$ જો અને માત્ર જો $\mSat{M^*}{!A}[{[w]}]$,
કારણ કે જૂનું નિદર્શન~$!A$નું વિરોધી નિદર્શન હોય તો નવું પણ
એવું જ જોઈએ. આ માટે વિભાજન તથા~$\mModel{M^*}$નો પ્રાપ્યતા
સંબંધ યોગ્ય રીતે વ્યાખ્યાયિત કરવો પડશે.

આ રીત કેવી રીતે ચાલે તે જોવા પહેલાં કલ્પો કે પ્રાપ્યતા સંબંધ
સર્વત્ર છે. $\mSat{M}{\Box !B}[w]$ જો અને માત્ર જો કોઈ
$v \in W$ માટે $\mSat{M}{\Box !B}[v]$;~$\mModel{M^*}$ માટે પણ
એ જ વાત, માત્ર $[w]$ અને $[v]$ સાથે.
\sourcecorrection{OLMD-042}{મૂળમાં અહીં પ્રાપ્યતા સંબંધ જ ન હોય
એમ કહે છે, પરંતુ આગળ બૉક્સના કિસ્સામાં દરેક જગતના સૂત્ર-સત્યની
શરત વાપરે છે. તે શરત માટે અહીં સર્વત્ર પ્રાપ્યતાનો કિસ્સો
સ્પષ્ટ કર્યો છે.}
પ્રથમ વિચાર તરીકે, બે જગતો $u$ અને $v$ને સમકક્ષ (એ જ વર્ગનાં)
કહીએ જો તેઓ~$\mModel{M}$નાં બધાં !!{propositional variable}s
પર સહમત હોય, એટલે કે $\mSat{M}{p}[u]$ જો અને માત્ર જો
$\mSat{M}{p}[v]$. $V^*(p) = \Setabs{[w]}{\mSat{M}{p}[w]}$ લો.
આપણું લક્ષ્ય $\mSat{M}{!A}[w]$ જો અને માત્ર જો
$\mSat{M^*}{!A}[{[w]}]$ બતાવવાનું છે. સ્પષ્ટ છે કે આ અનુમાનથી
સાબિત કરીશું: આધાર-કિસ્સો $!A \equiv p$ હશે. પહેલાં ધારો કે
$\mSat{M}{p}[w]$. ત્યારે વ્યાખ્યા મુજબ $[w] \in V^*(p)$,
તેથી $\mSat{M^*}{p}[{[w]}]$.
\sourcecorrection{OLMD-043}{મૂળના આ આધાર-કિસ્સામાં નિયુક્તિના
ચિહ્ન પછી વિધાનચલનો આર્ગ્યુમેન્ટ છૂટી ગયો છે; અગાઉની વ્યાખ્યા
સાથે તેને મેળ ખાતો કર્યો છે.}
હવે ધારો કે $\mSat{M^*}{p}[{[w]}]$. તેનો અર્થ
$[w] \in V^*(p)$, એટલે કે~$w$ને સમકક્ષ કોઈ $v$ માટે
$\mSat{M}{p}[v]$. પરંતુ ``$w$ એ $v$ને સમકક્ષ છે'' એટલે
``$w$ અને $v$ એ બધા સરખા !!{propositional variable}sને સત્ય
બનાવે છે''; તેથી $\mSat{M}{p}[w]$. હવે અનુમાનનું પગલું લો,
દાખલા તરીકે $!A \ident \lnot !B$. ત્યારે
$\mSat{M}{\lnot !B}[w]$ જો અને માત્ર જો $\mSat/{M}{!B}[w]$,
જો અને માત્ર જો $\mSat/{M^*}{!B}[{[w]}]$ (અનુમાનની
પૂર્વધારણાથી), જો અને માત્ર જો
$\mSat{M^*}{\lnot !B}[{[w]}]$. બીજા બિનમોડલ સંચાલકો માટે પણ
એ જ રીતે. $\Box$ માટે પણ ચાલે છે: ધારો કે
$\mSat{M^*}{\Box !B}[{[w]}]$. તેનો અર્થ કે દરેક $[u]$ માટે
$\mSat{M^*}{!B}[{[u]}]$. અનુમાનની પૂર્વધારણાથી દરેક $u$ માટે
$\mSat{M}{!B}[u]$. તેથી $\mSat{M}{\Box !B}[w]$.

સામાન્ય કિસ્સામાં~$\mModel{M^*}$નો પ્રાપ્યતા સંબંધ પણ વ્યાખ્યાયિત
કરવો પડે છે અને વાત વધુ જટિલ બને છે. ઉપરનું અનુમાન-પ્રમાણ ચાલે
તે માટે જરૂરી શરતો તેનો પ્રાપ્યતા સંબંધ~$R^*$ પૂરી કરતો હોય,
તો નિદર્શન~$\mModel{M^*}$ને \emph{ફિલ્ટ્રેશન} કહીશું.
ત્યારે કોઈપણ ફિલ્ટ્રેશન~$\mModel{M^*}$માં $[w]$ પર $!A$
સત્ય છે જો અને માત્ર જો~$\mModel{M}$માં~$w$ પર $!A$ સત્ય છે.
પરંતુ હવે ફિલ્ટ્રેશનો \emph{હોય છે} એ પણ બતાવવું પડે, એટલે કે
જરૂરી શરતો પૂરી કરે તેવો~$R^*$ વ્યાખ્યાયિત કરી શકાય છે.
આ શક્ય બને તે માટે જગતો $u$, $v$ બધાં !!{propositional variable}s
પર જ નહિ, પણ~$!A$નાં બધાં ઉપ-!!{formula}s પર સહમત હોય ત્યારે જ
સમકક્ષ ગણવા જોઈએ. $!A$નાં ઉપ-!!{formula}s સાન્ત સંખ્યામાં હોવાથી
ફિલ્ટ્રેશન સાન્ત છે તેની ખાતરી હજી મળે છે. ફિલ્ટ્રેશનને વ્યાખ્યાયિત
કરવાનો એક જ રસ્તો નથી; અને ફિલ્ટ્રેશનનો પ્રાપ્યતા સંબંધ જરૂરી
ગુણધર્મો (જેમ કે સ્વપ્રતિબિંબિત, પરંપરિત વગેરે) ધરાવે તે માટે
~$R^*$ની વ્યાખ્યામાં સર્જનાત્મકતા વાપરવી પડે છે.


\end{document}
